\section{Lean definitions and a path to complete formalization}
\label{sec:lean}

\subsection{The intended endpoint}

The final formal statement must construct a derivation in the embedded
arithmetic. In the repository's notation the target has the form
\[
                 \Gamma\ \vdash[\texttt{PRAAx}]\ \varphi,
\]
where $\varphi$ is a \texttt{Fm PTm n} and $\Gamma$ is a list of formulas
with the same numerical free variables. The actual type is
\texttt{Prf PRAAx}. A host Lean proof of a fact about natural numbers,
lists, rational numbers, or \texttt{CategoryTheory.Category} is a different
kind of theorem. It may guide the arithmetic proof, but is not its
replacement.

The inspected source defines \texttt{PR k} by constructors for zero,
successor, projection, composition, and primitive recursion. It defines
\texttt{PTm n} by variables and applications of these symbols. The
\texttt{PRAAx.ind} constructor requires a proof that its induction formula
is quantifier-free. The derived rule \texttt{PRA.qfInd} exposes this
restriction directly \cite{repo-pr}. These are the interfaces used here.

The repository pins Lean \texttt{4.32.0} and Mathlib \texttt{v4.32.0}; the
inspected manifest resolves Mathlib to commit
\texttt{81a5d257c8e410db227a6665ed08f64fea08e997}. These are reproducibility
targets, not a claim that this packet was compiled with them. No Lean
executable was available in the inspected environment. The new Lean files
are therefore explicitly \emph{uncompiled source}.

\subsection{Delivered definitions}

\begin{longtable}{@{}P{.32\linewidth}P{.61\linewidth}@{}}
\toprule
File & Mathematical content and status\\
\midrule\endhead
\path{lean/ObjectCategory.lean} & Source-level proofs with conclusions in
the actual \texttt{Prf PRAAx} calculus: binary congruence, inverse uniqueness,
an implication discharging its five hypotheses, both Yoneda determining
computations, and implication composition. Uncompiled.\\
\path{lean/FiniteCategoryTargets.lean} & Concrete finite-list data for maps,
categories, functors, diagrams, cones, cocones, presheaves and representable
transformations; exact propositions for limit and colimit producers and
Yoneda determination. Targets are definitions of propositions, not proofs.\\
\path{lean/MacaulayTargets.lean} & Concrete exponent and monomial enumeration,
Macaulay matrices, closure acceptance, action identities, evaluation vectors,
denominator grids, and complete-list targets over rational multivariate
polynomials. Uncompiled host specification.\\
\path{lean/CategoryBlueprint.md} & Data and proof details for contract maps,
presented setoid categories, Yoneda, finite universal constructions, and the
object-theory bridge.\\
\path{lean/FORMALIZATION.md} & Ordered implementation milestones, exact
proof obligations, coding bounds and invariants, build instructions, and
status distinctions.\\
\bottomrule
\end{longtable}

The finite-category specification deliberately uses \texttt{List Nat} for
table storage. For example, composition is lookup at index $g\,a+f$ in an
$a^2$-entry table, where $a$ is the arrow count. Invalid list lookup returns
zero, while validity and typing are explicit predicates. The specification
does not introduce an opaque arbitrary-function evaluator.

The Macaulay specification uses host \texttt{MvPolynomial (Fin r) Rat}
for readability. The human proof specifies the eventual numerical encoding:
finite lists of exponent tuples and rational coefficients, with equality
checked coefficientwise. \texttt{ClosureAccepted} is an existential
linear-algebra \emph{specification}. A first formalization stage must implement
the bounded row solver and prove its Boolean acceptance equivalent to that
specification. Treating the existential formula itself as a Boolean algorithm
would leave the main production obligation unproved.

\subsection{One actual object-proof pattern}

For a fixed binary PR symbol \texttt{comp}, let \texttt{bin comp g f}
denote the arithmetic term for $g\circ f$. The delivered lemma
\texttt{inverse\_unique} consumes object derivations of
\begin{align*}
 g\circ1_B&=g, & f\circ h&=1_B,\\
 g\circ(f\circ h)&=(g\circ f)\circ h, & g\circ f&=1_A,\\
 1_A\circ h&=h.
\end{align*}
It returns an object derivation of $g=h$, by the exact five-step equality
chain already proved. The companion \texttt{inverse\_unique\_imp} introduces
the conjunction of those five assertions as an antecedent and proves the
implication from the empty context. Thus the categorical equations are
\emph{hypotheses in the theorem}, not extra axioms added to \texttt{PRAAx}.
An application derives them from the chosen category's validity predicate.

The delivered \texttt{yoneda\_component} similarly consumes a naturality
instance at $1_A$ and the equation $1_A\circ f=f$. Two uses of equality
reasoning yield
\[
 \alpha_X(f)=F(f)(\alpha_A(1_A)).
\]
This is the determining step, not an assertion that the whole Yoneda
construction has already been compiled. Its converse, typing, and global
table properties remain explicit subsequent lemmas.

These source expressions use numerical equality. For a genuine setoid
presentation, replace each equality by the appropriate relation formula and
supply its symmetry, transitivity, and operation congruence proofs. The same
human derivation then compiles through those rules. No quotient-equality law
is obtained merely by reusing arithmetic equality substitution.

\subsection{Implementation sequence for the main lemmas}

\paragraph{Stage 1: numerical primitives.}
Implement list coding, lookup, update, finite conjunction, first-success
search, rank/select in a retained list, bounded products and sums, signed
integers, rational arithmetic, and exponent tuples as \texttt{PR} syntax.
For each primitive, prove its defining equations and a small set of lookup or
algebraic laws in \texttt{Prf PRAAx}. Record the numerical validity guards.
No general soundness theorem is needed.

\paragraph{Stage 2: Boolean predicates.}
Implement category, functor, naturality, cone, finite cover, overlap, matrix,
and polynomial-identity tests. Prove their equivalence to the displayed
bounded formulas using the prefix proofs in Lemma~\ref{lem:booleanization}.
This stage is what permits the object proof to obtain a particular category
law from an accepted table, rather than simply assuming all laws.

\paragraph{Stage 3: fixed diagram proofs.}
First compile \texttt{ObjectCategory.lean} against the pinned repository and
audit the dependencies of its theorem declarations. Then instantiate its
local law hypotheses using Stage 2. Add the equally short split-monic,
terminal-uniqueness, naturality-composition, and adjunction calculations.
These are finite first-order derivations and do not need new induction.

\paragraph{Stage 4: uniform finite constructors.}
For limits, prove the filtered tuple list is exactly the compatible tuples,
then construct the rank mediator and prove uniqueness entrywise. For colimits,
prove bounded reachability and the least-representative quotient, then descend
cocones. For exponentials, prove digit summation and reconstruction, treating
empty cases before division. Finite Kan extensions compose these constructions
with the explicit formulas in Section~\ref{sec:category-framework}; no separate
existence principle is needed. For finite-contract arrows, prove membership in
the composite union and use the two input proof terms by implication cuts.

\paragraph{Stage 5: algebraic certificate synthesis.}
Prove the row solver correct using the invariant
$[R\mid w]=U[A\mid v]$, the pivot pattern, and an inverse-operation log.
Preserve the labels of the generator-multiple columns, so a solution vector
constructs the explicit ideal-membership polynomials. Prove coefficientwise
verification of each identity. Evaluate the identities at an arbitrary common
zero, prove primitive scaling and the denominator lemma, and obtain the finite
coordinate grid. Finally prove that exact filtering is sound and complete.
Each induction concerns a bounded table predicate; the candidate point stays
a free parameter until universal generalization.

\paragraph{Stage 6: gluing and coherence.}
Finite sheaf gluing is a first-success scan followed by a table loop. CRT uses
the finite extended-Euclid invariant and rational-free integer calculations.
For surjective towers, fix the supplied numerical initial value, implement the
bounded lift scan, and use a QF validity invariant. The negative computability
results are metatheorems about proposed uniform procedures; they should not be
misrepresented as unproved object-theory axioms.

\paragraph{Stage 7: native transport.}
Only after obtaining the intended arithmetic derivation invoke the repository's
\texttt{pra\_iff\_native}, supplied by the module
\path{ConciseIPRA.NativeBiInterpretation}
\cite{repo-native}. The theorem relates the actual ordinary-iPRA proof
judgement to the native translation of its context and conclusion. It changes
the presentation of the proof; it does not eliminate a mathematical hypothesis,
create a missing bound, or turn an arbitrary host theorem into an iPRA theorem.

\subsection{A compact invariant inventory}

\begin{longtable}{@{}P{.24\linewidth}P{.32\linewidth}P{.35\linewidth}@{}}
\toprule
Construction & Finite resource & Induction content\\
\midrule\endhead
Contract composition & Lengths of the stored query lists & Membership in the
flattened union; local implication for each demanded question.\\
Category/functor tables & Object, arrow and tuple counts & Retained prefix
contains exactly the candidates passing the bounded test.\\
Finite limit & $\prod_i d_i$ candidates & Soundness and completeness of
compatible retained tuples; rank/select equations.\\
Finite colimit & $N=\sum_i d_i$ rounds; temporary paths of length at most $2N$
& Reachability with intermediate vertices below the current round.\\
Macaulay elimination & Matrix dimensions; finite row-operation log & Tracked
matrix identity, inverse operations, and pivot pattern.\\
Coordinate grid & $\prod_k(2R_k+1)$ candidates & Every retained point is a zero;
every zero at an already scanned grid index was retained.\\
Finite skeleton & Length of the supplied cover & Retained entries are
inequivalent and cover the scanned prefix.\\
Surjective tower & Stage index and given finite next-stage bound & The current
state is valid; the selected next state maps to it.\\
\bottomrule
\end{longtable}

For a loop bounded by a final parameter, guard its invariant by being within
that bound. For matrix or table identities, bounded quantifiers are removed
through their Boolean representatives before \texttt{PRA.qfInd} is used.
This inventory gives a finite completion route; it is not a claim that the
implementation has already been carried out.
